% part: intuitionistic-logic
% chapter: propositions-as-types
% section: proof-terms

\documentclass[../../../include/open-logic-section]{subfiles}

\begin{document}

\olfileid{int}{pty}{ter}

\olsection{Terma Bukti}

Kita menehi label variabel (contone, $!A^x$) marang asumsi ing \Log{N1i}
minangka label pelepasan. Ing \Log{N2i}, saben sekuen $\Gamma \Sequent !A$
nduweni dhaftar pasangan variabel--!!{formula} ing konteks~$\Gamma$ ing
sisih kiwa. Mula, ing \Log{N2i}, saben sekuen ing !!a{proof} ora mung
ngandhut informasi ngenani apa sing !!{proof} kasebut wis !!{prove}
nganti titik kasebut (yaiku $!A$), nanging uga asumsi endi sing isih
mbukak ing saben titik (yaiku $\Gamma$). Kepiye yen kita uga nyakup
informasi \emph{kepiye} $!A$ !!{prove} ing saben sekuen? Kanggo iki,
kita bisa netepake sistem \emph{mawa label terma} \Log{tN2i}, kanthi
sekuen awujud $\Gamma \Sequent N : !A$. $\Gamma$~lan $!A$ padha karo
sadurunge: himpunan pasangan awujud $x: !B$ lan !!{formula} sing wis
dituduhake dening sub-!!{proof} sing dipungkasi ing sekuen kasebut
minangka akibat saka asumsi ing~$\Gamma$. $N$~bakal dadi sawijining
terma sing ngode !!{proof} sing dipungkasi ing sekuen kasebut.

Terma sing kita pasang ing saben sekuen digawe saka variabel $x$, $y$,
\dots, kanthi simbol fungsi sing ngode inferensi sing wis ditindakake.
Kita mbutuhake simbol fungsi dhewe kanggo saben aturan inferensi.
Simbol fungsi sing cocog karo aturan !!{introduction} diarani
``konstruktor'', lan sing cocog karo aturan !!{elimination} diarani
``destruktor''. Saben simbol fungsi nampa argumen sing cacahe padha
karo cacah premis aturan kasebut. Contone, simbol fungsi kanggo
\Intro{\land} lan \Elim{\land}[1] bisa awujud $c^\land$ (rong panggonan)
lan $d^\land_1$ (siji panggonan). Ing \Log{tN2i}, aturan iki dadi
\[
\Axiom$\Gamma_1 \fCenter N: !A$
\Axiom$\Gamma_2 \fCenter M: !B$
\RightLabel{\Intro{\land}}
\BinaryInf$\Gamma_1, \Gamma_2 \fCenter c^\land(N, M) : !A \land !B $
\DisplayProof
\quad
\Axiom$\Gamma \fCenter N: !A \land !B$
\RightLabel{\Elim{\land}[1]}
\UnaryInf$\Gamma \fCenter d^\land_1(N) : !A$
\DisplayProof
\]
Yen sawijining aturan inferensi ngeculake asumsi, yaiku nduweni premis
sing nyebut !!{formula} mawa label~$x:!A$ sing ora ana ing konteks
konklusine, kita uga kudu menehi tandha iki ing terma bukti. Simbol
fungsi sing digandhengake karo aturan kaya mangkono \emph{ngiket}
variabel sing cocog. Contone, aturan \Intro{\lif} ngowahi premis
$x:!A, \Gamma \Sequent !B$ dadi $\Gamma \Sequent !A \lif !B$.
Konstruktor $c^{\lif}$ nampa siji argumen, nanging uga kudu menehi
tandha yen variabel $x$ dadi kaiket. Iki carane terma bukti nyatakake
yen asumsi mawa label~$x$ wis diuculake. Contone, yen terma bukti
kanggo premis iku~$N$, kita bisa nulis terma bukti kanggo konklusi
minangka $c^{\lif}([x]N)$, lan aturane minangka
\[
\Axiom$x:!A, \Gamma \fCenter N: !B$
\RightLabel{\Intro{\lif}}
\UnaryInf$\Gamma \fCenter c^{\lif}([x]N) : !A \lif !B$
\DisplayProof
\]

Supaya bisa ngrakit maneh !!a{proof} kanthi jangkep, kita mbutuhake
informasi tambahan. Yen !!{formula}s premis sawijining aturan ora
nemtokake !!{formula} konklusine kanthi jangkep, kita kudu nambahake
informasi kasebut ing terma. Iki dumadi ing \Intro{\lif}, mula kita
bakal nulis $c^{\lif}([x:!A]N)$. Iki uga dumadi ing \Intro{\lor}
lan~\FalseInt, mula kita nggunakake simbol fungsi sing ngandhut
informasi kasebut, contone,
\[
\Axiom$\Gamma \fCenter N:!A$
\RightLabel{\Intro{\lor}[1]}
\UnaryInf$\Gamma \fCenter c^{\lor{!B}}_1(N):!A \lor !B$
\DisplayProof
\quad
\Axiom$\Gamma \fCenter N:\lfalse$
\RightLabel{\FalseInt}
\UnaryInf$\Gamma \fCenter d^\lfalse_{!A}(N):!A$
\DisplayProof
\]

Kanthi gagasan iki kita bisa ngrakit sistem dedhuksi alami kanthi gaya
sekuen, kanthi saben sekuen awujud $\Gamma \Sequent N : !A$, lan
$N$~minangka \emph{terma bukti}.

Ana sesambungan raket antarane terma bukti kasebut lan terma ing
\emph{kalkulus lambda mawa tipe}. Kanggo njaga sesambungan iki, kita
bakal nggunakake simbol sing digunakake ing literatur kalkulus lambda,
dudu simbol konstruktor lan destruktor umum ing ndhuwur. Contone,
tinimbang $c^{\lif}([x:!A]N)$ kita nulis $\lambd[\typeof{x}{!A}][N]$;
tinimbang $d^{\lif}(N,M)$ kita nulis $(NM)$;
$\pair{N}{M}$ ngganti $c^\land(N,M)$, lan $\proj{i}{N}$ ngganti
$d^\land_i(N)$.

\begin{defn}
  Terma bukti diasilake kanthi induktif miturut aturan ing ngisor iki:
  \begin{enumerate}
  \item Saben variabel $x$ iku terma bukti (Aksioma).
  \item Yen $x$ iku variabel, $!A$ iku !!a{formula}, lan $N$ iku terma
    bukti, $\lambd[\typeof{x}{!A}][N]$ uga terma bukti~(\Intro\lif).
  \item Yen $N$ lan $M$ iku terma bukti, $(NM)$ uga terma
    bukti~(\Elim\lif).
  \item Yen $N$ lan $M$ iku terma bukti, $\pair{N}{M}$ iku terma
    bukti~(\Intro\land).
  \item Yen $N$ iku terma bukti, $\proj{i}{N}$ uga terma bukti,
    kanthi $i$ minangka $1$ utawa~$2$~(\Elim\land[i]).
  \item Yen $N$ iku terma bukti, lan $!A$ iku formula,
    $\inj[!A]{i}{N}$ iku terma bukti, kanthi $i$ minangka $1$
    utawa~$2$~(\Intro\lor[i]).
  \item Yen $M$, $N_1$, $N_2$ iku terma bukti, lan $x_1$, $x_2$ iku
    variabel, $\dcase{M}{x_1}{N_1}{x_2}{N_2}$ iku terma
    bukti~(\Elim\lor).
  \item Yen $N$ iku terma bukti lan $!A$ iku !!a{formula},
    $\abort{!A}{N}$ iku terma bukti~(\FalseInt).
  \end{enumerate}
\end{defn}

Ora saben terma bukti iku terma bukti saka !!a{proof}. Contone,
terma $\pair{N}{M}$ ngode bukti sing dipungkasi kanthi inferensi
\Intro{\land}, mula !!{formula} pungkasane iku konjungsi~$!A \land !B$.
Iki ora bisa dadi premis utama aturan \Elim{\lif}, mula
$(\pair{N}{M}P)$ ora bisa dadi terma bukti saka bukti sing bener.
Mung terma bukti sing manut aturan \Log{tN2ip} sing dadi terma bukti
\emph{bener}. Aturan kasebut diwenehake ing \olref{tab:tN2ip}.

\subfile{rules-tN2}

\end{document}
